% ===================================================================
%  TABLO No 5 — Ev ve yapılışı.
%  Traduction 2026 d'apres texto/io, controlee sur texto/fr et
%  texto/hi. Consigne : voir texto/tr/10-tablo-01.tex.
%
%  Le tableau 5 est un DIALOGUE, et les deux volumes composent le nom
%  du parleur en petites capitales : on garde \textsc{}, que le rendu
%  laisse tomber en gardant le texte.
%
%  Les renvois du plan sont des LETTRES — (a) le couloir, (b) le
%  salon... — et non des numeros. On les ecrit comme l'ido les ecrit,
%  y compris la ou l'imprimeur a compose « (1) » pour « (l) ».
%
%  LES DEUX NOMS PROPRES SE TRANSCRIVENT, ILS NE SE TRADUISENT PAS.
%  Le hindi rend « Leblanc » par le mot qui veut dire « blanc » ; le
%  turc ne le peut pas, parce qu'il ecrit deja tous les autres noms
%  du livret a l'oreille. « Bay Löblan » et « Bay Ru » suivent donc
%  la meme regle que Jan et Anri au tableau 1 : une seule regle pour
%  toute la colonne, ou aucune.
% ===================================================================

%%K t05-apar-1 apar
\VUsaut{6.0mm}
\VUtitre{15.0pt}{60}{TABLO No 5}
\VUsaut{1.4mm}
\VUtitre{11.4pt}{0}{Ev ve yapılışı.}
\VUsaut{2.2mm}

%%K t05-tit-1 sub
\VUpk{10.2pt}{40}{İyi bir karşılaşma.}
\VUsaut{3.0mm}

%%K t05-01-1 p
\textsc{Jan}. --- Amca, bir \VUgras{evin} nasıl yapıldığını pek bilmek
isterim.

%%K t05-01-2 p
\textsc{Amca} \textsuperscript{(80)}. --- Kolay, oğlum; benimle gel de sana
Bay Bernar'ın \textsuperscript{(79)} yaptırdığı ve benim oturmaya gideceğim
evi göstereyim.

%%K t05-01-3 p
\textsc{Jan}. --- Peki bu evin \VUgras{planını} \textsuperscript{(76)} kim
çizdi?

%%K t05-01-4 p
\textsc{Amca}. --- \VUgras{Mimar} \textsuperscript{(75)}; pek temiz bir çizim
yaptı, o da kabul edildi, ve \VUgras{müteahhit} \textsuperscript{(77)} işe
başladı.

%%K t05-01-5 p
\textsc{Jan}. --- Peki müteahhit kimdir?

%%K t05-01-6 p
\textsc{Amca}. --- Bay Löblan, arkadaşın Jak'ın babası. Mimar Bay Ru ile
birlikte işçilere yol gösterir.

%%K t05-01-7 p
Bak! Bay Löblan tam zamanında geliyor, \VUgras{metresiyle}
\textsuperscript{(78)}, Bay Ru da planıyla.

%%K t05-01-8 p
Günaydın efendim. Nasılsınız?

%%K t05-01-9 p
\textsc{Bay Ru}. --- Pek iyiyim, teşekkür ederim. Günaydın Jan; oğlan ne kadar
büyümüş!

%%K t05-01-10 p
\textsc{Amca}. --- Gerçekten, tam küçük bir adam oldu. Pek ağırbaşlıdır;
gördüğü her şeyi ona açıklamak gerekir. Bugün bir evin nasıl yapıldığını
öğrenmek istiyor.

%%K t05-tit-2 sub
\VUpk{10.2pt}{40}{İşçiler.}
\VUsaut{3.0mm}

%%K t05-01-11 p
\textsc{Bay Löblan}. --- Öyleyse benimle gel, küçük dostum, sana hemen
anlatayım, çünkü öğrenmek isteyen çocuklardan pek hoşlanırım.

%%K t05-01-12 p
Elinde \VUgras{küreği} \textsuperscript{(82)} olan şu işçiyi görüyor musun :
o \VUgras{kazıcıdır} \textsuperscript{(81)}. Su borusunu döşemek için bir
\VUgras{hendek} \textsuperscript{(46)} kazıyor. Duvarcının dört
\VUgras{duvarı} kaldırabilmesi için evin durduğu toprağı da o kazmıştı. O
duvarların toprak içinde kalan bölümüne \VUgras{temel} \textsuperscript{(2)}
denir. Temelin çevirdiği yer \VUgras{bodrumu} \textsuperscript{(3)} yapar. Bu
bodrumda \VUgras{havalandırma penceresi} \textsuperscript{(5)} denen küçük bir
pencere, bir \VUgras{kapı} \textsuperscript{(4)} ve bir merdiven vardır.

%%K t05-01-13 p
\VUgras{Duvarcı} \textsuperscript{(108)} duvarları kaldıra kaldıra çıktı,
\VUgras{kapılara} \textsuperscript{(8)} ve \VUgras{pencerelere}
\textsuperscript{(17)} yer bırakarak.

%%K t05-01-14 p
\textsc{Jan}. --- Ama ben bu duvarcıyı hiç göremiyorum!

%%K t05-01-15 p
\textsc{Bay Löblan}. --- Onun evdeki işi artık bitmiştir. Şurada
\VUgras{ahırın} \textsuperscript{(54)} yanındadır, \VUgras{iskelesine}
\textsuperscript{(110)} çıkmış, \VUgras{çamaşırhanenin}
\textsuperscript{(56)} duvarını örüyor.

%%K t05-01-16 p
Yanında bir mala, bir tekne ve bir \VUgras{çekül} \textsuperscript{(109)}
vardır. Ama bu adları aklında tutamazsın.

%%K t05-01-17 p
\textsc{Jan}. --- Neden tutamayayım, amcamdan küçük bir mala ile bir tekne
isteyeceğim, çeküle gelince onu bir sicim parçasıyla kendim yaparım.

%%K t05-tit-3 sub
\VUsaut{3.0mm}
\VUpk{10.2pt}{40}{Gereçler.}
\VUsaut{3.0mm}

%%K t05-01-18 p
\textsc{Amca}. --- Aferin, pek güzel. Ama söyle bakalım Jan, bir ev yapmak
için ne gerekir?

%%K t05-01-19 p
\textsc{Jan}. --- \VUgras{Yontma taş} \textsuperscript{(59)} ile
\VUgras{harç} \textsuperscript{(65)} gerekir.

%%K t05-01-20 p
\textsc{Amca}. --- Doğru. Şu büyük şapkalı adama bak, \VUgras{arabasının}
\textsuperscript{(136)} üstünde. O \VUgras{arabacıdır}
\textsuperscript{(135)}. Her gün buraya \VUgras{taş blokları}
\textsuperscript{(60)} getirir. Arabasını avluda boşaltır, ve
\VUgras{taşçılar} \textsuperscript{(83)} blokları yontar ve
\VUgras{taş testereleriyle} \textsuperscript{(85)} keser. Bir \VUgras{işçi}
\textsuperscript{(146)} onları duvarcıya ulaştırır, o da yerlerine oturtur.

%%K t05-01-21 p
Bu arada \VUgras{harççı} \textsuperscript{(87)} harcı hazırlar. Ona
\VUgras{kum} \textsuperscript{(62)}, \VUgras{kireç} \textsuperscript{(63)} ve
\VUgras{su} \textsuperscript{(64)} gerekir. Kireç yanında bir \VUgras{çuval}
\textsuperscript{(71)} içindedir; bir \VUgras{duvarcı yardımcısı}
\textsuperscript{(111)} ona \VUgras{el arabasıyla} \textsuperscript{(112)} kum
getirir, ve \VUgras{çırak} \textsuperscript{(113)} tulumbadadır
\textsuperscript{(58)}. Bir \VUgras{kova} \textsuperscript{(114)} dolduruyor.
Harç birazdan hazır olacak. \VUgras{Harç taşıyıcısı} \textsuperscript{(107)}
onu alıp merdivenden iskeleye çıkarır; orada bir duvarcı evin o kanadını
bitiriyor.

%%K t05-01-22 p
Şimdi söyle, \VUgras{çatıyı} \textsuperscript{(31)} yapan işçiye ne denir?

%%K t05-01-23 p
\textsc{Jan}. --- O \VUgras{çatıcıdır} \textsuperscript{(118)}.

%%K t05-tit-4 sub
\VUsaut{3.0mm}
\VUpk{10.2pt}{40}{Evin çatısı.}
\VUsaut{3.0mm}

%%K t05-01-24 p
\textsc{Bay Löblan}. --- Evet, ama önce \VUgras{dülger}
\textsuperscript{(115)} gelir.

%%K t05-01-25 p
Dülger \VUgras{çatı iskeletini} \textsuperscript{(35)} yapar.
\VUgras{Kirişleri} \textsuperscript{(37)} \VUgras{takozlarla}
\textsuperscript{(117)} birleştirir. Yukarıdaki, geniş kadife pantolonlu o
işçiler dülgerdir. Biri küçük bir kirişi tutuyor, öteki de
\VUgras{baltasıyla} \textsuperscript{(116)} bir takoz yontuyor.
\VUgras{Bacanın} \textsuperscript{(41)} yanındaki çatıcıdır.
\VUgras{Merteklere} (*) çıtaları çakar, çıtalara da \VUgras{arduvazları}
\textsuperscript{(66)}. Kimi zaman kiremit döşer.

%%K t05-noto-1 noto
\VUnotes{143.2mm}{%
(*) \VUgras{Balk-o} = almanca \textit{Balken}, ingilizce \textit{joist},
fransızca \textit{solive}, italyanca \textit{travicello}, rusça
\textit{balka}, ispanyolca ile portekizce \textit{viga}. Daha benimsenmemiş
teknik bir kök, ama burada fransızca \textit{solive} sözünün karşılığı olsun
diye önerilir; o ne \textit{trabo}\,dur (fransızca \textit{poutre}) ne de
küçük bir kiriş. Bir döşemeyi ve bir tavanı taşıyan, dört köşe yontulmuş
tahtadır.%
}

%%K t05-01-26 p
Ahırın çatısı \VUgras{kiremitten} \textsuperscript{(55)}, evinki
\VUgras{arduvazdan} \textsuperscript{(32)}, verandanınki de
\VUgras{çinkodandır} \textsuperscript{(24)}.

%%K t05-01-27 p
\textsc{Jan}. --- Çinko çatıları da çatıcı mı yapar?

%%K t05-01-28 p
\textsc{Bay Löblan}. --- Hayır, o \VUgras{çinkocunun}
\textsuperscript{(121)} ya da \VUgras{tesisatçının} işidir --- evin çatısına
bir \VUgras{çatı penceresi} \textsuperscript{(38)} oturtan da odur.
\VUgras{Çatı oluklarını} \textsuperscript{(43)} ve \VUgras{iniş borularını}
\textsuperscript{(42)} da o takar. Çinko ile tenekeyi lehimler.
\VUgras{Yıldırımsavarı} \textsuperscript{(40)} ile \VUgras{rüzgârgülünü}
\textsuperscript{(39)} \VUgras{mahyaya} \textsuperscript{(36)} çakan da odur.

%%K t05-01-29 p
\textsc{Jan}. --- Demek ev bitti?

%%K t05-tit-5 sub
\VUsaut{3.0mm}
\VUpk{10.2pt}{40}{Aletler.}
\VUsaut{3.0mm}

%%K t05-01-30 p
\textsc{Bay Löblan}. --- Bitmedi. \VUgras{Sıvacı} \textsuperscript{(99)}
\VUgras{tuğlayla} \textsuperscript{(61)} evi ayrı ayrı odalara bölen duvarları
yapar. \VUgras{Sıvayı} \textsuperscript{(70)} \VUgras{tavanlara}
\textsuperscript{(26)} ve duvarlara çeker. Şu sırada \VUgras{verandanın}
\textsuperscript{(33)} tavanında çalışıyor. Duvarcı gibi onun da bir
\VUgras{malası} \textsuperscript{(100)} ile bir \VUgras{teknesi}
\textsuperscript{(101)} vardır.

%%K t05-01-31 p
Sonra dülger kapıları, pencereleri, \VUgras{ahşap döşemeleri}
\textsuperscript{(16)} ve \VUgras{panjurları} \textsuperscript{(19)} yapar.

%%K t05-01-32 p
Şu sırada avluda \VUgras{tezgâhının} \textsuperscript{(90)} başında
çalışıyor. Bir \VUgras{testeresi} \textsuperscript{(94)} vardır,
\VUgras{tahtayı} \textsuperscript{(69)} kesmek için, bir \VUgras{rendesi}
\textsuperscript{(91)}, bir \VUgras{mengenesi} \textsuperscript{(92)}, bir
\VUgras{keskisi} \textsuperscript{(94 \textit{bis})}, bir \VUgras{pergeli}
\textsuperscript{(95)} ve tahtadan bir \VUgras{tokmağı}
\textsuperscript{(95 \textit{bis})}.

%%K t05-01-33 p
\textsc{Jan}. --- Vay, öyleyse dülger olmak duvarcı olmaktan da eğlencelidir.
Amca, bana bir tezgâh, bir testere ve bir rende alır mısınız? Medor'a bir ev
yaparım.

%%K t05-01-34 p
\textsc{Amca}. --- Olur oğlum.

%%K t05-01-35 p
\textsc{Jan}. --- Ne sevindim! Peki sokak kapısının yanında duran şu adam
kimdir?

%%K t05-tit-6 sub
\VUsaut{3.0mm}
\VUpk{10.2pt}{40}{Boyama.}
\VUsaut{3.0mm}

%%K t05-01-36 p
\textsc{Bay Löblan}. --- O \VUgras{boyacıdır} \textsuperscript{(102)}.
Merdiveninin üstünde durur, bir kap \VUgras{boya} \textsuperscript{(103)} ile
\VUgras{fırçasını} \textsuperscript{(104)} tutarak. Daha uzun dayansın diye
sokak kapısının tahtasını boyuyor.

%%K t05-01-37 p
Odalardaki duvar kâğıdını da o yapıştırır.

%%K t05-01-38 p
\VUgras{Döşemeci} \textsuperscript{(107)}, \VUgras{merdivene}
\textsuperscript{(106)} çıkmış \VUgras{balkondadır} \textsuperscript{(28)},
bir \VUgras{perde} \textsuperscript{(29)} takıyor.

%%K t05-01-39 p
Büyük pencerenin yanında duran işçi \VUgras{camcıdır}
\textsuperscript{(96)}. \VUgras{Camları} \textsuperscript{(18)}
\VUgras{macunla} \textsuperscript{(73)} tutturur. Bodrumun kapısı yanında diz
çökmüş öteki işçi \VUgras{çilingirdir} \textsuperscript{(97)}. Bir
\VUgras{kilit} \textsuperscript{(6)} takıyor. \VUgras{Eğesi}
\textsuperscript{(98)} yanında durur.

%%K t05-01-40 p
\textsc{Jan}. --- \VUgras{Çekiciyle} \textsuperscript{(84)} bir demir parçasına
vuran adam da mı çilingirdir?

%%K t05-01-41 p
\textsc{Bay Löblan}. --- Doğrusunu söylemek gerekirse o \VUgras{demircidir}
\textsuperscript{(125)}. \VUgras{Demir işlerini} \textsuperscript{(67)}
dövüyor, \VUgras{elektrikçi} \textsuperscript{(124)} için, ki
\VUgras{araba barakasının} \textsuperscript{(51)} çatısındadır. Bir \VUgras{ocağı}
\textsuperscript{(126)}, bir \VUgras{örsü} \textsuperscript{(127)} ve bir
\VUgras{kerpeteni} \textsuperscript{(128)} vardır.

%%K t05-01-42 p
Elektrikçi telefonun \VUgras{bakır telini} \textsuperscript{(44)} bağlıyor.

%%K t05-tit-7 sub
\VUpk{10.2pt}{40}{Bahçıvanlık.}
\VUsaut{3.0mm}

%%K t05-01-43 p
\textsc{Amca}. --- \VUgras{Avlunun} \textsuperscript{(45)} dibinde,
\VUgras{parmaklık} \textsuperscript{(47)} ile \VUgras{araba kapısı}
\textsuperscript{(48)} yanında, \VUgras{bahçıvanı} \textsuperscript{(129)}
görüyorsun. Küçük \VUgras{bahçeyi} \textsuperscript{(49)} hazırlıyor.
\VUgras{Beliyle} \textsuperscript{(131)} toprağı kazmıştır, tohum ekiyor ve
\VUgras{tırmıkla} \textsuperscript{(130)} düzlüyor. Sonra
\VUgras{süzgeçli kovasıyla} \textsuperscript{(132)} \VUgras{tarhları}
\textsuperscript{(150)} sulayacak ve bitkiler bitecek.

%%K t05-01-44 p
Atı gözetip yemini vermiş olan \VUgras{seyis} \textsuperscript{(133)},
\VUgras{faytonu} \textsuperscript{(52)} temizliyor --- sünger, bir
\VUgras{el tulumbası} \textsuperscript{(134)} ve bir kova suyla. Sonra onu
araba barakasına koyacak ve ağır \VUgras{sürgüyle} \textsuperscript{(53)}
kapıyı kapatacak.

%%K t05-01-45 p
Hadi bakalım, artık yetinmişsindir; anlatmak yeter de arttı.

%%K t05-01-46 p
\textsc{Jan}. --- Evet amca, ama evin içini de görmek isterim.

%%K t05-01-47 p
\textsc{Amca}. --- O da yarına. Bay Ru sana planı açıklar ve her odanın adını
söyler.

%%K t05-tit-8 sub
\VUsaut{3.0mm}
\VUpk{10.2pt}{40}{Evin iç düzeni.}
\VUsaut{2.4mm}

%%K t05-apar-2 apar
{\centering\textit{(Plana bakınız.)}\par}
\VUsaut{2.4mm}

%%K t05-01-48 p
\textsc{Mimar}. --- Gel Jan, sana evin ayrı ayrı bölümlerini göstereyim.

%%K t05-01-49 p
Hadi \VUgras{zemin kata} \textsuperscript{(7)} girelim. İşte \VUgras{zilin}
düğmesi \textsuperscript{(11)}. Üç \VUgras{basamağı}
\textsuperscript{(14)} çıkarız, ki \VUgras{giriş merdiveninindir}
\textsuperscript{(9)}, sonra \VUgras{eşiği} \textsuperscript{(10)} geçeriz,
ki \VUgras{sokak kapısınındır} \textsuperscript{(8)}, ve işte \VUgras{merdivenin}
\textsuperscript{(13)} altındayız.

%%K t05-01-50 p
\textsc{Jan}. --- Burası koridordur.

%%K t05-01-51 p
\textsc{Mimar}. --- Hayır, burası \VUgras{antredir} \textsuperscript{(12)}.
\VUgras{Koridor} \textsuperscript{(a)} diptedir. Sağda işte \VUgras{salon}
\textsuperscript{(b)} ile \VUgras{yemek odası} \textsuperscript{(c)}; yemek
odasının karşısında \VUgras{mutfak} \textsuperscript{(d)} ile \VUgras{kiler}
\textsuperscript{(e)}; sonra ön yana doğru \VUgras{çalışma odası}
\textsuperscript{(f)}. \VUgras{Veranda} \textsuperscript{(23)},
\VUgras{camlı kapısıyla} \textsuperscript{(25)} birlikte, salonun yanındadır.

%%K t05-01-52 p
Şimdi merdiveni çıkacağız. \VUgras{Tırabzana} \textsuperscript{(15)} tutun ve
basamakları say. On dört tanedir. İşte \VUgras{sahanlık}
\textsuperscript{(m)} : birinci \VUgras{kattayız} \textsuperscript{(27)}.

%%K t05-tit-9 sub
\VUsaut{3.0mm}
\VUpk{10.2pt}{40}{Birinci kat.}
\VUsaut{3.0mm}

%%K t05-01-53 p
Merdivenin karşısında \VUgras{helâ} \textsuperscript{(20)} ile
\VUgras{giyinme odası} \textsuperscript{(j)} vardır. Sağda güzel bir
\VUgras{yatak odası} \textsuperscript{(g)}, ki iki penceresi evin
\VUgras{cephesine} \textsuperscript{(1)} açılır. Solda \VUgras{banyo}
\textsuperscript{(1)} ile \VUgras{konuk odası} \textsuperscript{(i)} vardır,
ki içinde büyük bir \VUgras{niş} \textsuperscript{(h)} bulunur. Yatak odasının
yanında \VUgras{çocuk odası} \textsuperscript{(k)} vardır. Geniş bir
\VUgras{terasa} \textsuperscript{(21)} açılır. Bu terasın altında bir helâ daha
\textsuperscript{(20)} vardır, ve duvara bitişik işte \VUgras{çan}
\textsuperscript{(22)}, ki yemek için çalınır.

%%K t05-01-54 p
\textsc{Jan}. --- Hadi \VUgras{ikinci kata} çıkalım.

%%K t05-01-55 p
\textsc{Mimar}. --- İkinci kat yoktur. Çatının altında yalnız
\VUgras{tavan araları} \textsuperscript{(33)} ile depo \textsuperscript{(34)}
vardır. Turumuz bitti, artık inebiliriz.

%%K t05-01-56 p
\textsc{Jan}. --- Teşekkür ederim efendim, pek ilgi çekiciydi. Gördüğüm her
şeyi babama anlatacağım.
\VUsaut{4.0mm}
\VUfilet{22mm}
